Desde la perspectiva de la ingeniería de sonido y la producción musical, el proceso de mezcla representa una de las etapas más complejas y determinantes dentro de la creación de una obra sonora. En particular, la mezcla vocal requiere una serie de decisiones técnicas relacionadas con ecualización, control dinámico, espacialidad y procesamiento armónico, las cuales influyen directamente en la claridad, presencia e integración de la voz dentro del arreglo musical [1]. Sin embargo, en la práctica profesional se reconoce que no existe una única forma correcta de realizar una mezcla, ya que cada ingeniero o productor desarrolla su propio enfoque técnico y estético dependiendo del género musical, las características de la grabación y la intención artística del proyecto. [2] Esta diversidad de criterios genera un escenario en el que los productores con menor experiencia pueden enfrentar dificultades al momento de estructurar cadenas de procesamiento vocal coherentes.

Desde un punto de vista práctico, esta situación es especialmente relevante en el contexto actual de la producción musical digital y los entornos de home studio, donde una gran cantidad de productores independientes asumen simultáneamente los roles de grabación, edición y mezcla [1]. Aunque existen numerosas herramientas digitales y plugins de procesamiento de audio, estas no siempre ofrecen orientación clara sobre cómo organizar o estructurar una cadena de mezcla adecuada para una voz específica [3]. Como consecuencia, muchos productores trabajan mediante ensayo y error, lo que puede prolongar los tiempos de producción y generar resultados inconsistentes entre diferentes proyectos.

En este contexto, el desarrollo de una plataforma inteligente capaz de generar propuestas de cadenas de mezcla vocal puede representar un aporte significativo para el campo de la producción musical [4]. En lugar de reemplazar el criterio creativo del ingeniero de sonido, una herramienta de este tipo podría funcionar como un sistema de asistencia que ofrezca configuraciones iniciales de procesamiento basadas en características acústicas de la voz [5]. De esta manera, el productor o ingeniero mantendría el control sobre las decisiones finales de mezcla, pero contaría con una guía técnica que facilite la organización de su flujo de trabajo y reduzca la incertidumbre al momento de construir cadenas de procesamiento vocal. Este enfoque busca fortalecer la relación entre tecnología e ingeniería de sonido, promoviendo el uso de herramientas inteligentes como apoyo para la toma de decisiones dentro del proceso creativo de la producción musical [1].
